\documentclass[12pt]{article}
\input{iati_structure.tex}
\renewcommand{\bottomtitlespace}{2cm}

\usepackage[english]{babel}

\begin{document}

\renewcommand{\headerLang}{English}
\renewcommand{\problemName}{AB23. TRIANGLE}

\renewcommand{\headerLeft}{IATI Day 2, Senior group}
\renewcommand{\headerRightFirst}{XVII INTERNATIONAL ADVANCED TOURNAMENT IN INFORMATICS}
\renewcommand{\headerRightSecond}{ONLINE 2026}

\renewcommand{\tl}{$10$ sec.}
\renewcommand{\ml}{$1024$ MB}
\problem{Problem \problemName}
\addauthor{Author: Boris Mihov}{2026}{04}{19}

Նոր ազգային կոմիտեի ղեկավարությունը, ներկայացված A և B խմբերի համակարգողների կողմից, ցանկանում է պահպանել տրված խնդիրների խիստ սիլաբուսը՝ ազգային հավաքականի ընտրության համար։ Այդ նպատակով նրանք պետք է ներկայացնեն երկրաչափության խնդիր։ Ավելին — մինչ այժմ տրված ինտերակտիվ խնդիրների քանակը չի համապատասխանում պահանջվող մակարդակին։ \texttt{IOI}-ի ընտրության ճգնաժամը լուծելու համար Sashka-ն որոշեց տալ խնդիր, որը միաժամանակ ինտերակտիվ է և երկրաչափական։ Այն ունի հետևյալ ձևակերպումը․

Ժյուրին թաքցրել է ${1,2,3,...,N}$ բազմության մի $P_0,P_1,...,P_{N-1}$ տեղափոխություն։ Դուք պետք է գտնեք այդ տեղափոխությունը։ Դրա համար կարող եք ժյուրիին տալ հետևյալ հարցը․ «Հնարավո՞ր է կազմել դրական մակերես ունեցող եռանկյուն $P_A,P_B,P_C$ կողմերով»։

Նկատենք, որ հայտնի է (եռանկյան անհավասարությունից), որ դա հնարավոր է այն և միայն այն դեպքում, երբ՝

\[P_A + P_B > P_C,\ P_A + P_C > P_B \quad \text{and} \quad P_B + P_C > P_A.\]

Գրեք ծրագիր \textbf{\texttt{triangle}}, որը պարունակում է \texttt{solve} ֆունկցիա, որը կկոմպիլացվի ժյուրիի ծրագրի հետ միասին և կհաղորդակցվի դրա հետ՝ տալով վերը նկարագրված տեսակի հարցեր։ Իր աշխատանքի ավարտին այն պետք է որոշի տեղափոխությունը։

\subsection{Իրականացման մանրամասներ}
Դուք պետք է իրականացնեք \texttt{solve} ֆունկցիան:
\begin{lstlisting}
 std::vector<int> solve(int N)
\end{lstlisting}
\begin{itemize}
	\item $N$: տեղափոխության երկարությունը։
\end{itemize}

Այս ֆունկցիան կկանչվի յուրաքանչյուր թեստի համար $T$ անգամ — յուրաքանչյուր ենթաթեստի համար մեկ անգամ, որտեղ բոլոր ենթաթեստերում $N$-ը նույնն է, և այն պետք է վերադարձնի տվյալ ենթաթեստի համար թաքցված տեղափոխությունը։ Դրա համար ձեր ծրագիրը կարող է կանչել ժյուրիի \texttt{query} ֆունկցիան։

\begin{lstlisting}
 bool query(int A, int B, int C)
\end{lstlisting}
\begin{itemize}
	\item $A$, $B$, $C$: ինդեքսներն են $P_A, P_B, P_C$ կողմերի։
\end{itemize}

Ֆունկցիան վերադարձնում է \texttt{true}, եթե $P_A,P_B,P_C$ կողմերով եռանկյուն գոյություն ունի, և \texttt{false}՝ հակառակ դեպքում։

\subsection{Սահմանափակումներ}
\vspace{0.1em}
\begin{itemize}
	\item $N = T = 1~000$
\end{itemize}

\subsection{Subtask}
\begin{table}[H]
	\begin{tblr}{width=\linewidth,colspec={|Q[c,m]|Q[c,m]|X[c,m]|}}
		\hline
		\textbf{Subtask} & \textbf{Points} & \textbf{Description}\\
		\hline
		$1$ & $100$ & Ենթախնդիրը բաղկացած է $1$ թեստից, որում կա $T=1~000$ պատահական և միատեսակ գեներացված տեղափոխություններ, յուրաքանչյուրը կազմված $N=1~000$ տարրերից։  \\
		\hline
	\end{tblr}
	\caption*{Տրված ենթախնդրի համար միավորները տրվում են միայն այն դեպքում, եթե բոլոր թեստերն անցնում են։}
\end{table}
\FloatBarrier

\subsection{Գնահատումը}

Թող $Q_{contestant}$-ը լինի ձեր ծրագրի կողմից կատարված հարցումների միջին քանակը՝ մեկ \texttt{solve} կանչի համար տվյալ մեկ թեստում, և թող $Q_{target}=8770$։ Այդ դեպքում ձեր միավորը այս խնդրի համար կլինի․

\[
\begin{cases}
0 & \text{if } Q_{contestant} > 2\times 10^6, \\
100 & \text{if } Q_{contestant} \le Q_{target}, \\
100 \times \max\left(0.15,\; 1-\sqrt{1-\left(\frac{Q_{target}}{Q_{contestant}}\right)^{0.55}}\right) & \text{if } Q_{contestant} > Q_{target}.
\end{cases}
\]

\subsection{Փոխգործակցման օրինակ}
For this sample interaction $N=4$, $T=2$.
\begin{table}[H]
	\begin{tblr}{|c|c|}
		\hline
		\textbf{Contestant action} & \textbf{Jury action} \\
		\hline
		& \texttt{solve(4)} \\
		\hline
		\texttt{query(0, 0, 0)} & \texttt{true} \\
		\hline
        \texttt{query(0, 1, 2)} & \texttt{false} \\
		\hline
        \texttt{query(0, 1, 3)} & \texttt{false} \\
		\hline
        \texttt{query(0, 2, 3)} & \texttt{true} \\
		\hline
        \texttt{query(1, 2, 3)} & \texttt{false} \\
		\hline
        \texttt{return \{3, 1, 2, 4\};} & \\
        \hline
        & \texttt{solve(4)} \\
		\hline
		\texttt{query(0, 1, 2)} & \texttt{false} \\
		\hline
        \texttt{query(0, 1, 3)} & \texttt{true} \\
		\hline
        \texttt{query(0, 2, 3)} & \texttt{false} \\
		\hline
        \texttt{query(1, 2, 3)} & \texttt{false} \\
		\hline
        \texttt{return \{4, 2, 1, 3\};} & \\
		\hline
	\end{tblr}
\end{table}
\FloatBarrier

\subsection{Լոկալ թեստավորում}

Մուտքային տվյալների ձևաչափ՝
\begin{itemize}
	\item տող $1$: երեք ամբողջ թիվ՝ $T$, $N$ և $R$ -- տեղափոխությունների չափը, ենթաթեստերի քանակը և կատարման ռեժիմը։
	\item եթե $R = 1$, մուտքի հաջորդ տողը պարունակում է մեկ թիվ՝ $S$, որը կօգտագործվի որպես պատահական գեներատորի seed, և լոկալ գրեյդերը կստեղծի պատահական և միատեսակ տեղափոխություններ։
	\item եթե $R = 2$, մուտքը շարունակվում է հետևյալ կերպ՝
	\item տողեր $2$-ից մինչև $(1+T)$: $\{1,2,\dots,N\}$ բազմության տեղափոխություններ։
\end{itemize}

Ելքային ձևաչափ՝
\begin{itemize}
	\item տող $1$: սխալի հաղորդագրություն կամ հարցումների միջին քանակը ենթաթեստերի համար, եթե բոլոր տեղափոխությունները ճիշտ են հայտնաբերվել։
\end{itemize}
	
\end{document}